\documentclass[10pt,a4paper]{amsart}
\usepackage[margin=19mm]{geometry}
\usepackage{amsmath,amssymb,mathtools}
\usepackage[scr=rsfs,cal=euler]{mathalfa}
\usepackage{xfrac}
\usepackage[unicode,hidelinks,bookmarks=false]{hyperref}
\newcommand{\ie}{i.e.\ }
\newcommand{\cf}{cf.\ }
\newcommand{\resp}{respectively\ }
\newcommand{\Subsection}{\subsection}
\providecommand{\nobreakdash}{\nobreak\hbox{-}}
\setlength{\emergencystretch}{2em}
\multlinegap0pt
\usepackage{amssymb}
\newtheorem{corollary}{Corollary}
\title{On the Minimality of the Conductor for Elliptic Curve $L$-Functions}
\author{K. Lakshmanan}
\date{Source: arXiv:2506.20175v2 (CC BY 4.0)}
\begin{document}
\maketitle

\noindent\emph{Source and licence.} On the Minimality of the Conductor for Elliptic Curve $L$-Functions. Version arXiv:2506.20175v2 (CC BY 4.0), distributed by arXiv under the Creative Commons Attribution 4.0 International licence, http://creativecommons.org/licenses/by/4.0/, as recorded in the arXiv licence metadata for this exact version and shown on the abstract page https://arxiv.org/abs/2506.20175v2. Full licence notice: \texttt{licenses/arxiv-2506.20175-CC-BY-4.0.txt}.\par\smallskip

\noindent\textit{Editorial scope.} Counted unit: the article's corollary corollary (source0.tex lines 239--251). The article's complete original proof of it is delivered with it (source0.tex lines 253--261, one \texttt{proof} environment of 9 lines). No mathematics is written, repaired or re-derived: the statement and the proof are the article's own lines, in the article's own notation, and no retained line is substituted or re-flowed.  No further source-local context is needed: every cross-reference of every retained line resolves inside the two delivered spans.  Nothing was omitted from any retained span, so the package carries no omission record.  No retained line contains a citation, so the wrapper carries no bibliography and no bibliography entry.  No retained line contains a figure environment, an image-inclusion command or drawing code, so no image is delivered and the package declares no asset dependency.  Editorial formatting only, in addition: an AMS \texttt{amsart} preamble that restates the article's own declarations and macros used by the retained lines, namely the environments \texttt{corollary} on separate counters, together with the standard packages those lines need.  The retained environments print in this document as Corollary 1 in order of appearance, whereas the article prints the same items with the numbering of its own full document.  The complete native source of the article as distributed in the arXiv e-print for this exact version is bundled unchanged as \texttt{sources/180475/source0.tex}.
	\begin{corollary}[Conditional Rank Unboundedness]
		Suppose there exists an infinite family of elliptic curves \( \{E_n\} \) over \( \mathbb{Q} \), and an arithmetic invariant \( \Phi(E) \) such that:
		\begin{enumerate}
			\item \( \Phi(E_n) < N_{E_n} \) for all \( n \),
			\item \( \Phi(E_n) \to \infty \) as \( n \to \infty \),
			\item \( \operatorname{rank}(E_n) \ll \log \Phi(E_n) \) holds for all \( n \).
		\end{enumerate}
		Then the ranks \( \operatorname{rank}(E_n) \) are unbounded:
		\[
		\sup_n \operatorname{rank}(E_n) = \infty.
		\]
		But this contradicts the main theorem, unless \( \Phi(E_n) \ge N_{E_n} \) for all sufficiently large \( n \).
	\end{corollary}

	\begin{proof}
		If such a family \( \{E_n\} \) and function \( \Phi(E) \) exist, then the analytic rank is controlled by a sequence of invariants growing strictly slower than the conductors \( N_{E_n} \). Since \( \log \Phi(E_n) < \log N_{E_n} \), this would give a tighter rank bound than the classical inequality:
		\[
		\operatorname{rank}(E_n) \ll \log N_{E_n}.
		\]
		But the main theorem shows that any such functional equation or bound must involve the full conductor \( N_{E_n} \), and no strictly smaller invariant \( \Phi(E) \) can appear in the analytic structure of the \( L \)-function. Hence, the assumption leads to a contradiction unless \( \Phi(E_n) \ge N_{E_n} \) for all large \( n \), i.e., unless \( \Phi \) essentially coincides with the conductor.
		
		Thus, if \( \Phi(E_n) < N_{E_n} \) and grows unbounded, then \( \operatorname{rank}(E_n) \to \infty \), and rank must be unbounded.
	\end{proof}

\end{document}
